% ===== PRÉAMBULE CANONIQUE — à coller VERBATIM en tête de CHAQUE .tex =====
\documentclass[11pt,a4paper]{article}
\usepackage[utf8]{inputenc}\usepackage[T1]{fontenc}\usepackage{lmodern}
\usepackage[french]{babel}
\usepackage{amsmath,amssymb,mathtools}\usepackage{icomma}
\usepackage[version=4]{mhchem}          % \ce{...} : équations & formules
\usepackage{siunitx}\sisetup{locale=FR} % \qty{}{}, \si{}, \ang{}
\usepackage{chemfig}                    % \chemfig{...} : structures développées
\usepackage{xcolor}\usepackage{colortbl}
\usepackage{tikz}\usepackage{pgfplots}\pgfplotsset{compat=1.18}
\usepackage[most]{tcolorbox}\usepackage{enumitem}\usepackage{array,booktabs}
\usepackage[a4paper,margin=2.2cm]{geometry}
\usetikzlibrary{arrows.meta,calc,angles,quotes,positioning,patterns,backgrounds,shapes.geometric}
\definecolor{primary}{RGB}{222,49,99}\definecolor{primarydark}{RGB}{150,22,55}
\definecolor{defcol}{RGB}{0,114,178}\definecolor{methcol}{RGB}{52,92,148}
\definecolor{excol}{RGB}{0,135,90}\definecolor{attncol}{RGB}{200,40,55}
\tcbset{boxbase/.style={enhanced,breakable,boxrule=0.9pt,arc=2.5pt,
  left=3mm,right=3mm,top=2.3mm,bottom=2.3mm,fonttitle=\bfseries,coltitle=white}}
% --- boîtes TUTORIEL (noms EXACTS, ne pas renommer) ---
\newtcolorbox{cle}[1][]{boxbase,colback=defcol!6,colframe=defcol,
  title={\textsc{À retenir}\ifx\relax#1\relax\relax\else\textmd{ : \itshape #1}\fi}}
\newtcolorbox{methode}[1][]{boxbase,colback=methcol!7,colframe=methcol,sharp corners,
  title={\textsc{Méthode}\ifx\relax#1\relax\relax\else\textmd{ --- \itshape #1}\fi}}
\newtcolorbox{exemplebox}[1][]{boxbase,colback=excol!5,colframe=excol,
  title={\textsc{Exemple}\ifx\relax#1\relax\relax\else\textmd{ : \itshape #1}\fi}}
\newtcolorbox{piege}[1][]{boxbase,colback=attncol!6,colframe=attncol,
  title={\textsc{Piège}\ifx\relax#1\relax\relax\else\textmd{ : \itshape #1}\fi}}
\newtcolorbox{remarque}[1][]{boxbase,colback=black!4,colframe=black!45,
  title={\textsc{Remarque}\ifx\relax#1\relax\relax\else\textmd{ : \itshape #1}\fi}}
% --- boîtes EXERCICE / CORRIGÉ (noms EXACTS, auto-comptées) ---
\newtcolorbox[auto counter]{exo}[1][]{boxbase,colback=primary!4,colframe=primary,
  title={\textsc{Exercice}~\thetcbcounter\ifx\relax#1\relax\relax\else\textmd{ --- \itshape #1}\fi}}
\newtcolorbox[auto counter]{corrige}[1][]{boxbase,colback=excol!4,colframe=excol,
  title={\textsc{Corrigé}~\thetcbcounter\ifx\relax#1\relax\relax\else\textmd{ --- \itshape #1}\fi}}
\newcommand{\R}{\mathbb{R}}\newcommand{\N}{\mathbb{N}}\newcommand{\Z}{\mathbb{Z}}\newcommand{\ds}{\displaystyle}
% (un \usepackage{fancyhdr}+en-tête est possible ; il est ignoré côté web)
% =========================================================================
\begin{document}

\begin{center}{\large\bfseries\color{primarydark} Thème 2 — Niveau Difficile}\end{center}
\emph{Données (\si{\gram\per\mole}) :} $\ce{H}=1$, $\ce{O}=16$, $\ce{Na}=23$, $\ce{S}=32$.
\quad $V_m = \SI{22.4}{\litre\per\mole}$ (CNTP), $N_{\mathrm{A}} = 6{,}02\times10^{23}\ \si{\per\mole}$.

\begin{exo}[identifier un gaz par sa masse molaire]
Un échantillon gazeux de masse $\SI{8.5}{\gram}$ occupe $\SI{5.6}{\litre}$ aux CNTP.
\begin{enumerate}[label=\alph*)]
  \item Quelle est sa quantité de matière~?
  \item En déduire sa masse molaire $M$.
  \item Sachant qu'il s'agit d'un hydrure de formule \ce{H2X}, identifie l'élément \ce{X}.
\end{enumerate}
\end{exo}

\begin{exo}[le piège des unités]
Un flacon vide pèse $\SI{13.45}{\gram}$. Rempli d'eau ($\rho = \SI{1.00}{\gram\per\milli\litre}$),
il pèse $\SI{21.45}{\gram}$. Rempli à ras bord d'un liquide inconnu, il pèse $\SI{57.89}{\gram}$.
\begin{enumerate}[label=\alph*)]
  \item Quel est le volume intérieur du flacon~?
  \item Quelle est la masse volumique du liquide inconnu, en \si{\gram\per\milli\litre}~?
  \item Exprime-la en \si{\gram\per\cubic\centi\metre}.
\end{enumerate}
\end{exo}

\begin{exo}[jongler avec les unités de masse]
L'or a une masse volumique $\rho = \SI{19.3}{\gram\per\cubic\centi\metre}$.
\begin{enumerate}[label=\alph*)]
  \item Quel volume (en \si{\cubic\centi\metre}) occupe $\SI{7.72}{\hecto\gram}$ d'or~?
        (\emph{rappel :} $\SI{1}{\hecto\gram} = \SI{100}{\gram}$.)
  \item Ce volume, exprimé en \si{\milli\litre}~?
\end{enumerate}
\end{exo}

\begin{exo}[du solide aux entités]
Un cube de soude \ce{NaOH} solide a une arête de $\SI{2}{\centi\metre}$ et une masse volumique
$\rho = \SI{2.0}{\gram\per\cubic\centi\metre}$. ($M(\ce{NaOH}) = \SI{40}{\gram\per\mole}$.)
\begin{enumerate}[label=\alph*)]
  \item Quelle est sa masse~?
  \item Quelle quantité de matière représente-t-il~?
  \item Combien de motifs \ce{NaOH} contient-il~?
\end{enumerate}
\end{exo}

\begin{exo}[comparer deux liquides]
À $\SI{4}{\celsius}$, on compare l'éthanol ($d = 0{,}79$, $M = \SI{46}{\gram\per\mole}$) et l'eau
($d = 1{,}00$, $M = \SI{18}{\gram\per\mole}$). On prélève $\SI{50}{\milli\litre}$ de chacun.
\begin{enumerate}[label=\alph*)]
  \item Lequel est le plus lourd~?
  \item Quelle quantité de matière contient chaque prélèvement~?
  \item Lequel contient le plus de molécules~? Commente.
\end{enumerate}
\end{exo}

\end{document}
